\section{四圈结果的独立外圈动量参考}
\label{sec:four-loop-independent}

This section concerns the fixed graph L4E3G205, with common squared mass
$M=1$ and symmetric Euclidean invariants $p_1^2=p_2^2=p_3^2=1$,
$p_1+p_2+p_3=0$.  Every original loop measure is $d^6k/\pi^3$; couplings,
vertex signs and diagram symmetry weights are excluded.  The reduction is
used only as an independent verification oracle.  Its implementation imports
neither the production Gaussian density and jets nor an existing
topology-specific integration kernel.  It does not replace graph-driven
production execution by a formula registry.

\subsection{The actual graph and the eight recursive terms}

The graph has nine vertices, twelve internal edges and external legs at
vertices $6,7,8$.  In original edge order its edges are
\begin{equation}
\begin{split}
 &(0,5),(0,5),(0,8),(1,4),(1,4),(1,7),\\
 &(2,3),(2,3),(2,7),(3,6),(4,5),(6,8).
\end{split}
\end{equation}
Its only proper divergent subgraphs are the three vertex-disjoint bubbles
on edge sets $\{0,1\}$, $\{3,4\}$ and $\{6,7\}$.  Directly from the
disjoint-spinney Bogoliubov recursion,
\begin{equation}
 \overline R_G I=
 (1-t_{\gamma_1}^{2})(1-t_{\gamma_2}^{2})(1-t_{\gamma_3}^{2})I,
 \qquad R_G I=(1-t_G^0)\overline R_G I.
 \label{eq:four-independent-eight}
\end{equation}
The product has eight terms.  Each local Taylor operator acts on the full
boundary momentum of its own bubble.  This is an elementary factorization
of this particular recursive expression, not an invocation of a general
forest or unverified project $Z$ identity.

Tracing the outer cycle gives
\begin{equation}
 6\longrightarrow\operatorname{bubble}(3,2)\longrightarrow7
 \longrightarrow\operatorname{bubble}(1,4)
 \longrightarrow\operatorname{bubble}(5,0)\longrightarrow8
 \longrightarrow6.
\end{equation}
The three side insertion counts are therefore $(1,2,0)$, rather than one
bubble on every side.  A side containing $n$ bubbles has $n+1$ ordinary
outer propagators, so define
\begin{align}
 \Sigma_R(s;M)
 &=\int_0^1\bigl[(M+s u(1-u))
       \log(1+s u(1-u)/M)-s u(1-u)\bigr],du,\\
 L_n(s;M)&=\frac{\Sigma_R(s;M)^n}{(s+M)^{n+1}}.
 \label{eq:four-independent-line}
\end{align}
The finite subtraction integrates to $-s/6$ and is retained throughout.
There is no extra factor $1/2$ for an insertion.  At finite common
original-edge Schwinger cutoff $\eta>0$, the independent inner bubble
integrations factor exactly in \eqref{eq:four-independent-eight}.

For the symmetric point we may use side counts $(0,1,2)$ and choose
$q_0=k$, $q_1=k+p_1$, $q_2=k-p_3$.  The complete overall degree-zero
subtraction is then
\begin{equation}
 J=\int\frac{d^6k}{\pi^3}
 \left[L_0(k^2)L_1((k+p_1)^2)L_2((k-p_3)^2)
       -\frac{\Sigma_R(k^2)^3}{(1+k^2)^6}\right].
 \label{eq:four-independent-outer}
\end{equation}
Writing three copies of $L_1$ would describe a different graph, although its
zero-point term would happen to be the same.  This makes the off-shell
branch mapping an essential part of the independent check.

\subsection{The three-dimensional integral and its stable evaluation}

Set $r=k^2$ and embed the external plane using
$p_1=(1,0)$ and $p_3=(-1/2,\sqrt3/2)$.  The projection of a uniform $S^5$
direction onto this plane has density
$2(1-x^2-y^2)/\pi$ on the unit disk.  With
$x=\sqrt v\cos\theta$, $y=\sqrt v\sin\theta$, the normalized measure and
shifted squared momenta are
\begin{align}
 \frac{d^6k}{\pi^3}
 &\longrightarrow\frac12 r^2,dr\,2(1-v)\,dv\,
                     \frac{d\theta}{2\pi},\\
 q_1^2&=r+1+2\sqrt r\,x,\qquad
 q_2^2=r+1+\sqrt r\,x-\sqrt{3r}\,y.
 \label{eq:four-independent-measure}
\end{align}
The directions $(x,y)$ and $(-x,-y)$ are paired before radial integration.
An exact independent closed form at $M=1$ is
\begin{equation}
 \Sigma_R(s)=\frac{(s+4)^{3/2}}{3\sqrt s}
       \operatorname{arsinh}(\sqrt s/2)-\frac43-\frac{4s}{9}.
 \label{eq:four-independent-bubble-closed}
\end{equation}
Near $s=0$, the implementation instead uses
\begin{equation}
 \Sigma_R(s)=\sum_{j=2}^{\infty}a_js^j,\qquad
 a_j=\frac{(-1)^j(j!)^2}{j(j-1)(2j+1)!},\quad
 a_2=\frac1{60},\quad
 \frac{a_{j+1}}{a_j}=-\frac{j-1}{2(2j+3)}.
\end{equation}
This closed form was also compared with a separate 70-digit
one-dimensional parameter integration at six scales; the largest observed
relative difference was $4.22\times10^{-15}$.

Let $\delta_\pm$ be the log ratio of the shifted product to its common
zero-point product, and set
$e=(\delta_++\delta_-)/2$, $o=(\delta_+-\delta_-)/2$.  The paired relative
difference is evaluated as
\begin{equation}
 \operatorname{expm1}(e)
       +e^e,2\sinh^2(o/2).
 \label{eq:four-independent-evenodd}
\end{equation}
At high radius the even and odd quantities are formed directly by analytic
Taylor derivatives and even/odd binomial grouping.  They are not obtained
by subtracting nearly equal floating logarithms.  The derivative
coefficients use 60--80-digit mpmath arithmetic before conversion to the
floating angular and radial quadrature; the full integral is not an
80-digit computation.

\subsection{A separate analytic radial-tail estimate}

For $r\geq4$, $|z|\leq1$ and $|a|=1$, write
$s=|k+za|^2$.  Then $r/4\leq s\leq9r/4$ and, with $\ell=\log r$,
the positive one-loop parameter formula gives
\begin{equation}
 0\leq\Sigma_R(s)\leq\frac{s\ell}{6},\qquad
 0\leq\Sigma_R'(s)\leq\frac\ell6,\qquad
 0\leq\Sigma_R''(s)\leq\frac1{6s}.
\end{equation}
Consequently, for $n=0,1,2$,
\begin{align}
 |L_n(s)|&\leq\frac{\ell^n}{6^n s},\\
 |L_n'(s)|&\leq\frac{(2n+1)\ell^n}{6^n s^2},\\
 |L_n''(s)|&\leq\frac{(4n^2+5n+2)\ell^n}{6^n s^3}.
 \label{eq:four-independent-line-bounds}
\end{align}
Put $g(z)=L_0(r)L_1(|k+zp_1|^2)L_2(|k-zp_3|^2)$.
Since $|s'|\leq3\sqrt r$ and $s''=2$, direct Leibniz differentiation gives
\begin{equation}
 |g''(z)|\leq\frac{20000}{27}\frac{\log^3r}{r^4}.
\end{equation}
The paired second-order Taylor remainder contributes a factor $1/2$;
the radial measure contributes another $1/2$.  Thus the whole omitted
outer tail above $T\geq4$ is bounded by
\begin{equation}
 |\operatorname{tail}(T)|\leq\frac{5000}{27T}
 \left(\log^3T+3\log^2T+6\log T+6\right).
 \label{eq:four-independent-tail}
\end{equation}
At $T=\operatorname{expm1}(48)$ this is
$3.109\times10^{-14}$, as a display approximation to the analytic bound.

Without pairing, the first Taylor derivative satisfies
$|g'|\leq(64/9)\log^3r/r^{7/2}$, which is also radially integrable.
For the common-cutoff bubble, its value and first two $s$ derivatives are
bounded by the same positive uncut quantities.  The cutoff outer
propagators have the corresponding uncut derivative majorants through
order two.  The above bounds therefore justify angular pairing and
dominated removal of $\eta$ in the reduced outer integral, identifying
\eqref{eq:four-independent-outer} with the same BPH regulated limit.
This argument is not, by itself, a general $L^1$ proof for the unreduced
24-dimensional combined momentum expression.

\subsection{A separate exact-coefficient Taylor truncation estimate}

Write $\Sigma_R(s)=s^2\phi(s)$ with
\begin{equation}
 \phi(s)=\int_0^1h(u)^2,du\int_0^1
             \frac{1-z}{1+s h(u)z},dz,\qquad h(u)=u(1-u).
\end{equation}
For $|w|=\varrho<1$, the ratio
$\phi(r(1+w))/\phi(r)$ is a positive-weight average of $(1+vw)^{-1}$,
$0\leq v\leq1$.  Its modulus is at most $(1-\varrho)^{-1}$ and its real
part is at least $(1-\varrho)/(1+\varrho)^2$.  It has a consistent analytic
logarithm, and
\begin{align}
 \left|\log\frac{\phi(r(1+w))}{\phi(r)}\right|
 &\leq2\log(1+\varrho)-\log(1-\varrho)+\frac\pi2,\\
 |\log(1+vw)|&\leq-\log(1-\varrho)+\frac\pi2.
\end{align}
Apply the two univariate Cauchy estimates separately to the shifted lines.
For counts $(0,1,2)$ there are three copies of $\log\Sigma_R$ and five
copies of $\log(1+s)$ in their combined log ratio.  At
$\varrho=9/10$ the sum of the boundary constants is less than $60$.

For a high-radius join $R$, Taylor order $J$, and
$q_0=(R^{-1}+2R^{-1/2})/\varrho<1$, put
$a=(J+1)/2$ and $D_*=60q_0/(1-q_0)$.  The Cauchy remainder, the exponential
mean-value estimate, and the radial prefactor bound
$\log^3r/432$ give
\begin{equation}
\begin{split}
 E_{\mathrm{Taylor}}\leq{}
 &\frac{e^{D_*}60q_0^{J+1}}{432(1-q_0)}\\
 &\times\left[
 \frac{\log^3R}{a}+\frac{3\log^2R}{a^2}
       +\frac{6\log R}{a^3}+\frac6{a^4}\right].
\end{split}
 \label{eq:four-independent-taylor-tail}
\end{equation}
For $R=1000,J=14$, this is $4.470\times10^{-15}$.
It bounds finite Taylor truncation with exact derivative coefficients.
It does not bound numerical coefficient rounding or the remaining
quadrature, and is distinct from \eqref{eq:four-independent-tail}.

\subsection{Calibration, actual refinements and preserved provenance}

The identical angular/radial scheme first reproduces the one-loop triangle
and the two-loop single-bubble BPH amplitude.  The latter's unit-point
reference is $-0.00617095361238479$.  In the last run its observed gap is
$1.27\times10^{-15}$; the one-loop calibration gap is
$6.25\times10^{-16}$.  The target four-loop QMC mean was not used to tune
the mapping, formula or refinement choices.

\begin{center}
\begin{tabular}{c c c}
\hline
Run & Radial/disk/half-angle & Four-loop value\\
\hline
Coarse & $12/20/24$ & $-0.0005158785412719462$\\
Fine & $20/32/40$ & $-0.0005158785412708389$\\
Finer & $28/48/64$ & $-0.0005158785412715456$\\
\hline
\end{tabular}
\end{center}
The first two runs use logarithmic panel width $2$, with respectively
$x=\log(1+r)\leq40,44$, join $R=10000$, Taylor orders $8,10$, and
60/70-digit derivative coefficients.  The last run uses panel width $1$,
$x\leq48$, join $R=1000$, Taylor order $14$ and 80-digit coefficients.
The last two results differ by $7.07\times10^{-16}$.  These are refinement
and calibration diagnostics, not a rigorous total quadrature enclosure.

The independent original script, three run JSON files, summary and
derivation notes are copied byte-for-byte into the documentation's
\texttt{automatic-four-loop-independent-evidence} directory, with original
times and hashes retained.  The portable tool
\path{tools/compare_automatic_four_loop_independent.py}
records its own hash and verifies the original script hash.  Its default
output is \path{results/automatic-four-loop-oracle-run.json}.
This reference is not bound to a transient production-source fingerprint;
a production comparison must preserve its own actual source, program,
execution and uncertainty records separately.
